\section*{第 \ref{ch:b2:magma-emplacement} 章——岩浆及其侵位}

\begin{solution}{exo:b2:magma-emplacement:1}
$T = \qty{1073}{K}$。当 $w = 1$ 时，$\ln w = 0$：$\log_{10}\eta = -3.545 +
9601/(1073 - 195.7) = -3.545 + 10.94 = 7.40$。当 $w = 4$ 时，$\ln 4 = 1.386$：
$A = -3.545 + 1.155 = -2.390$，$B = 9601 - 3283 = 6318$，$C = 195.7 + 44.7 =
240.4$，且 $\log_{10}\eta = -2.390 + 6318/832.6 = -2.390 + 7.588 = 5.20$。
粘度下降了 $10^{7.40 - 5.20} = 10^{2.20} \approx 160$ 倍。
\end{solution}

\begin{solution}{exo:b2:magma-emplacement:2}
流纹岩：$p_s = (6/0.44)^2 = \qty{186}{MPa}$，深度为 $\num{186e6}/(2600\times9.81)
= \qty{7.3}{km}$。玄武岩：$p_s = (0.5/0.34)^2 = \qty{2.16}{MPa}$，深度为
$\num{2.16e6}/(2600\times9.81) = \qty{85}{m}$。流纹岩在地壳深部即开始形成气泡，
而玄武岩仅在近地表才形成气泡（若不含 CO$_2$）。
\end{solution}

\begin{solution}{exo:b2:magma-emplacement:3}
$(1150 + 50)/2 = \qty{600}{\celsius}$ 与 $(1150 + 400)/2 = \qty{775}{\celsius}$。
冷岩石中的\omterm{def:b2:magma-emplacement:bodies}{岩墙}在与围岩接触处受到的冷却更强，
且其\omterm{def:b2:magma-emplacement:aureole}{冷凝边}更宽。
\end{solution}

\begin{solution}{exo:b2:magma-emplacement:4}
(a) \omterm{def:b2:magma-emplacement:bodies}{岩床}，与层理平行。(b) \omterm{def:b2:magma-emplacement:bodies}{岩墙}，切穿层理。
(c) \omterm{def:b2:magma-emplacement:bodies}{岩基}，面积超过 \qty{100}{km^2}，由多个侵入体组装而成。
(d) \omterm{def:b2:magma-emplacement:bodies}{岩盖}，它使顶板拱起。(e) \omterm{def:b2:magma-emplacement:bodies}{岩墙}：花岗岩无层状构造，
因此其中任何板状体都是\omterm{def:b2:magma-emplacement:bodies}{岩墙}，无论其产状如何。
\end{solution}

\begin{solution}{exo:b2:magma-emplacement:5}
$(1 - 1.35\phi)^{-2.5}$：取 $\phi = 0.05$，$0.9325^{-2.5} = 1.19$，故
\qty{119}{Pa.s}，而爱因斯坦公式 $1 + 2.5\times0.05 = 1.125$ 给出 \qty{113}{Pa.s}，
（两者接近，这在稀释极限下是应有的）。取 $0.20$：$0.73^{-2.5} =
2.20$，\qty{220}{Pa.s}；取 $0.40$：$0.46^{-2.5} = 6.97$，\qty{697}{Pa.s}；取 $0.60$：
$0.19^{-2.5} = 63.6$，约 \qty{6360}{Pa.s}。堵塞发生于 $\phi = 1/1.35 = 0.74$ 处。
\end{solution}

\begin{solution}{exo:b2:magma-emplacement:6}
$G = 300\times9.81 = \qty{2943}{Pa/m}$。(a) $\bar u = 2943\times1^2/(12\times100) =
\qty{2.45}{m/s}$；$q = \qty{2.45}{m^2/s}$；$Re = 2700\times2.45\times1/100 = 66$，层流；
$t = 10000/2.45 = \qty{4080}{s}$，约 \qty{1.1}{h}。(b) $\bar u =
2943/(12\times\num{e6}) = \qty{2.45e-4}{m/s}$，$t = \qty{4.08e7}{s} = 472$ 天；取
$w = \qty{10}{m}$ 时，$\bar u$ 为原来的 100 倍，即 \qty{0.0245}{m/s}，且 $t =
\qty{4.08e5}{s} = 4.7$ 天。(c) 宽 \qty{1}{m} 的\omterm{def:b2:magma-emplacement:bodies}{岩墙}约需
$0.5^2/\num{e-6} = \qty{2.5e5}{s}$，即 2.9 天冷却；宽 \qty{10}{m} 的\omterm{def:b2:magma-emplacement:bodies}{岩墙}则为 $5^2/\num{e-6} =
\qty{2.5e7}{s}$，即 289 天。玄武岩（\qty{1.1}{h}）及宽\omterm{def:b2:magma-emplacement:bodies}{岩墙}中的流纹岩
（4.7 天）均赶在\omterm{def:b2:magma-emplacement:bodies}{岩墙}凝固之前到达，而
\qty{1}{m} \omterm{def:b2:magma-emplacement:bodies}{岩墙}中的流纹岩（472 天）则不能。
\end{solution}

\begin{solution}{exo:b2:magma-emplacement:7}
每千克\omterm{def:b2:magma-emplacement:magma}{岩浆}中，水的物质的量为 $n = 10/18.015 = \qty{0.555}{mol}$，$nRT =
0.555\times8.314\times1173 = \qty{5414}{J}$；其余熔体占据的体积为 $0.99/2400 =
\qty{4.13e-4}{m^3}$。在 \qty{100}{MPa} 下：$V = \qty{5.41e-5}{m^3}$，分数为
$5.41/(5.41 + 41.3) = 0.12$。在 \qty{10}{MPa} 下：$V = \qty{5.41e-4}{m^3}$，分数为
$0.57$。在 \qty{0.1}{MPa} 下：$V = \qty{0.0541}{m^3}$，分数为 $0.992$。
\qty{100}{MPa} 处的数值最不可靠：该压强下的水蒸气密度大，
严重偏离理想气体（实际上大部分水仍会溶于熔体中）。
\end{solution}

\begin{solution}{exo:b2:magma-emplacement:8}
(a) 由 $2300 + 60z = 2650$ 得 $z = \qty{5.8}{km}$。(b) 熔体
密度为 \qty{2580}{kg/m^3}：$z = 280/60 = \qty{4.7}{km}$，约高 \qty{1.2}{km}。(c)
在 $z = \qty{7.8}{km}$ 处，岩石密度为 $2300 + 60\times7.8 = \qty{2770}{kg/m^3}$；
每立方米所受向上的作用力为 $(2770 - 2650)\times9.81 = \qty{1.2}{kN}$。
\end{solution}

\begin{solution}{exo:b2:magma-emplacement:9}
$v = 2\times300\times9.81\times(1000)^2/(9\times\num{e19}) = \qty{6.5e-11}{m/s}$，或
$\num{6.5e-11}\times\num{3.156e7} = \qty{2.1e-3}{m}$ 每年，即 \qty{2.1}{mm/\text{年}}。上升
\qty{10}{km} 需要 $\num{e4}/\num{2.1e-3} = \qty{4.8}{Ma}$。冷却时间尺度
为 $(1000)^2/\num{e-6} = \qty{e12}{s} = \qty{32}{ka}$：岩体在
上升约 $32000\times0.0021 \approx \qty{70}{m}$ 后凝固。\omterm{def:b2:magma-emplacement:magma}{岩浆}\omterm{def:b2:magma-emplacement:ascent}{底辟}无法穿越
这样的地壳；它需要温度更高、更柔软的外界。
\end{solution}

\begin{solution}{exo:b2:magma-emplacement:10}
(a) $x_{\mathrm{OH}} = \sqrt{K_2x_m} = \sqrt{K_1K_2f}$，故 $x = K_1f + \frac12\sqrt{K_1K_2f}$。
第一项与第二项之比为 $2\sqrt{K_1f/K_2}$，当
$f \to 0$ 时趋于零（平方根项占主导），随 $f$ 增大而无界增长（
线性项占主导）。(b) 由 $K_1 f = \frac12\sqrt{K_1K_2 f}$ 得 $K_1^2f^2 =
K_1K_2f/4$，故 $f = K_2/(4K_1)$。(c) $f = 0.2/(\num{4e-4}) = \qty{500}{MPa}$。在
\qty{100}{MPa} 下：$x_m = 0.0100$，$\frac12\sqrt{\num{e-4}\times0.2\times100} =
\frac12\times0.0447 = 0.0224$，故 $x = 0.0324$，羟基所占份额为
$0.0224/0.0324 = \qty{69}{\percent}$。
\end{solution}

\begin{solution}{exo:b2:magma-emplacement:11}
(a) $1000/\num{4e6} = \qty{2.5e-4}{km^3}$ 每年，即 $\qty{2.5e5}{m^3}$ 每年，
或 $\num{2.5e5}/\num{3.156e7} = \qty{0.0079}{m^3/s}$，约每秒 8 升。
(b) \omterm{def:b2:magma-emplacement:bodies}{岩墙}输送的流量为 $1000\times2\times0.8 = \qty{1600}{m^3/s}$，并
在 \qty{6.25e8}{s} 内输送 $\qty{e12}{m^3}$，约 20 年，仅为
\qty{4}{Ma} 的 $\num{5e-6}$。(c) 每个板状体容纳 $20\times\num{e4}\times\num{e4} = \qty{2e9}{m^3} =
\qty{2}{km^3}$，因此共 500 个板状体，平均每 $\num{4e6}/500 = 8000$ 年出现一个。
(d) $t = (10)^2/(4\times0.65^2\times\num{0.7e-6}) = \qty{8.5e7}{s}$，即 2.7 年。每个
板状体在下一个到来之前已凝固数千年：除非
\omterm{def:b2:magma-emplacement:magma}{岩浆}供给以短促的爆发式到来，否则\omterm{def:b2:magma-emplacement:bodies}{深成岩体}绝不会是一个大的\omterm{def:b2:magma-emplacement:magma}{岩浆}体，
这与大型\omterm{def:b2:magma-emplacement:bodies}{深成岩体}的定年结果相吻合。
\end{solution}

\begin{solution}{exo:b2:magma-emplacement:12}
(a) 固定 $x$，$\ln(T - T_0) = \text{常数} - \frac12\ln t - x^2/(4\kappa t)$，其
导数 $-1/(2t) + x^2/(4\kappa t^2)$ 在 $t^* = x^2/(2\kappa)$ 处为零。于是
$e^{-x^2/(4\kappa t^*)} = e^{-1/2}$，且 $\sqrt{4\pi\kappa t^*} = x\sqrt{2\pi}$，故 $\Delta
T_{\max} = (T_m - T_0)\,b\,e^{-1/2}/(x\sqrt{2\pi}) = (T_m - T_0)\,b/(x\sqrt{2\pi e})$。(b)
$\sqrt{2\pi e} = 4.13$：$x = 1100\times2/(100\times4.13) = \qty{5.3}{m}$，经过
$t^* = 5.3^2/(2\times\num{0.7e-6}) = \qty{2.0e7}{s}$ 后达到，约 230 天（仅在 2.7
倍\omterm{def:b2:magma-emplacement:bodies}{岩墙}宽度时，平面源近似较为粗糙）。(c) 加热作用随
$b/x$ 减小：薄\omterm{def:b2:magma-emplacement:bodies}{岩墙}携带的热量很少，岩石
只有在几个\omterm{def:b2:magma-emplacement:bodies}{岩墙}宽度之内才明显升温。
\end{solution}

\begin{solution}{pb:b2:magma-emplacement:1}
\textbf{1.} $\log_{10}\eta = -4.55 + 5510/(1423 - 611) = -4.55 + 6.786 = 2.24$;
$\eta = \qty{172}{Pa.s}$.

\textbf{2.} $-4.55 + 5510/762 = -4.55 + 7.231 = 2.68$，$\eta = \qty{480}{Pa.s}$，相差
$480/172 = 2.8$ 倍，对应 \qty{50}{K}。

\textbf{3.} $172\times(1 - 1.35\times0.20)^{-2.5} = 172\times0.73^{-2.5} = 172\times2.20 =
\qty{378}{Pa.s}$.

\textbf{4.} $\phi = 1/1.35 = 0.74$。野外证据表明，整体流动终止于
熔体分数为 \qtyrange{30}{50}{\percent} 处，即
晶体含量为 \qtyrange{50}{70}{\percent}，比此值偏早，因为实际晶体并非
等大的球体。

\textbf{5.} 基底：$2800 - 2710 = \qty{90}{kg/m^3}$，\omterm{def:b2:magma-emplacement:magma}{岩浆}较轻。
沉积层：$2450 - 2710 = \qty{-260}{kg/m^3}$，\omterm{def:b2:magma-emplacement:magma}{岩浆}较重。

\textbf{6.} $p = 2450\times9.81\times3000 = \qty{72.1}{MPa}$.

\textbf{7.} $90\times9.81 = \qty{883}{Pa/m}$.

\textbf{8.} $\bar u = 883\times2.0^2/(12\times378) = 3532/4536 = \qty{0.78}{m/s}$.

\textbf{9.} $q = 0.78\times2.0 = \qty{1.56}{m^2/s}$；沿 \qty{1.0}{km} 计，约
为 \qty{1560}{m^3/s}。

\textbf{10.} $Re = 2710\times0.78\times2.0/378 = 11$：层流。

\textbf{11.} $12000/0.78 = \qty{15400}{s}$，约 \qty{4.3}{h}。

\textbf{12.} $(1.0)^2/\num{0.7e-6} = \qty{1.43e6}{s}$，即 16.5 天，远长于
\qty{4.3}{h}：\omterm{def:b2:magma-emplacement:bodies}{岩墙}在\omterm{def:b2:magma-emplacement:magma}{岩浆}流动期间保持张开。

\textbf{13.} 在 \qty{3.0}{km} 处，\omterm{def:b2:magma-emplacement:magma}{岩浆}遇到比自身更轻的沉积层：
沉积层底部即其\omterm{def:b2:magma-emplacement:ascent}{中性浮力面}。在其上方，\omterm{def:b2:magma-emplacement:magma}{岩浆}密度将
大于外界，于是沿接触面侧向扩展，形成
\omterm{def:b2:magma-emplacement:bodies}{岩床}。

\textbf{14.} $p_s = (0.40/0.34)^2 = \qty{1.38}{MPa}$，该压强位于地表以下
$\num{1.38e6}/(2450\times9.81) = \qty{58}{m}$ 处。

\textbf{15.} $0.34\times\sqrt{72.1} = 0.34\times8.49 = \qty{2.9}{\percent}$，远高于
现有的 \qty{0.40}{\percent}，因此\omterm{def:b2:magma-emplacement:bodies}{岩床}\omterm{def:b2:magma-emplacement:magma}{岩浆}不出溶水。

\textbf{16.} CO$_2$ 的溶解度远低于水，它在
上升深部即饱和并形成气泡，因此玄武岩最初逸出的气体主要是 CO$_2$。

\textbf{17.} 每千克\omterm{def:b2:magma-emplacement:magma}{岩浆}，$n = 4.0/18.015 = \qty{0.222}{mol}$，$V = 0.222\times8.314\times1423/
\num{e5} = \qty{0.0263}{m^3}$；熔体占据 $0.996/2710 = \qty{3.68e-4}{m^3}$；气相所占
体积分数为 $0.0263/(0.0263 + 0.000368) = 0.986$，即 \qty{98.6}{\percent}。

\textbf{18.} $(1150 + 100)/2 = \qty{625}{\celsius}$.

\textbf{19.} 10 年后，$2\sqrt{\kappa t} = 2\sqrt{\num{0.7e-6}\times\num{3.156e8}} =
\qty{29.7}{m}$。该点距中部 $x = \qty{60}{m}$，$T = 100 +
525\,[\operatorname{erf}(-10/29.7) + \operatorname{erf}(110/29.7)] = 100 + 525\,(-0.365 +
1.000) = \qty{433}{\celsius}$.

\textbf{20.} 在中部，$T = 100 + 1050\operatorname{erf}(50/2\sqrt{\kappa t})$：
10 年后 $1050\operatorname{erf}(1.68) = 1050\times0.983$，$T = \qty{1132}{\celsius}$；经过
50 年后，$2\sqrt{\kappa t} = \qty{66.5}{m}$，$1050\operatorname{erf}(0.752) = 1050\times0.713$，
$T = \qty{848}{\celsius}$.

\textbf{21.} 在 $x = \qty{130}{m}$ 处：300 年后为 \qty{294}{\celsius}，
363 年后（最大值）为 \qty{295}{\celsius}，420
年后为 \qty{295}{\celsius}。若无潜热，\omterm{def:b2:magma-emplacement:bodies}{岩床}上方 \qty{80}{m} 处的岩石
始终未能达到 \qty{300}{\celsius}。

\textbf{22.} $\num{9e8}\times1.772/(\num{2.5e6}\times1050) = 0.608$.

\textbf{23.} $\lambda = 0.64$: $e^{-0.4096}/(0.64\times1.6346) = 0.6639/1.0461 = 0.635$;
$0.65$：$0.6554/1.0673 = 0.614$；$0.66$：$0.6469/1.0886 = 0.594$。数值 $0.608$
介于 $0.65$ 与 $0.66$ 之间，接近 $0.65$：$\lambda = 0.65$。

\textbf{24.} $100 + 1050/(1 + 0.642) = 100 + 640 = \qty{740}{\celsius}$。潜热
在凝固前沿释放，使\omterm{def:b2:magma-emplacement:magma}{岩浆}一侧保持在
\qty{1150}{\celsius} 的时间更长，因此接触带获得更多热量。

\textbf{25.} $(50)^2/\num{0.7e-6} = \qty{3.57e9}{s}$，113 年。

\textbf{26.} $t = (100)^2/(16\times0.65^2\times\num{0.7e-6}) = \num{e4}/\num{4.73e-6} =
\qty{2.11e9}{s}$，即 $\num{2.11e9}/\num{3.156e7}$：\omterm{def:b2:magma-emplacement:bodies}{岩床}在
\textbf{\qty[text-series-to-math=true]{67}{} 年}后完全凝固。这比
$(b/2)^2/\kappa$ 短，缩短因子为 $1/(4\lambda^2) = 0.59$：该时间尺度
仅决定数量级；\omterm{def:b2:magma-emplacement:bodies}{岩床}两侧壁都向冷的主岩导热，
且同时进行，速度足以带走潜热。
\end{solution}
